\section{Atomové jádro a elektronový obal}

\subsection{Atomové jádro}

\begin{itemize}
	\item Na stabilitu má vliv velikost vazebné energie jádra a poměr mezi počtem protonů a neutronů. U lehkých jader je poměr zhruba 1:1, se vzrůstajícím protonovým číslem dochází ke zvyšování přebytku neutronů.
	\item Vazebná energie je energie, která se uvolní při vzniku jádra z volných nukleonů
	\item Nejvíce stabilních jader má protonové i neutronové číslo sudé: \ce{$^{12}_{\ 6}$C}, \ce{$^{16}_{\ 8}$O}
	\item Naopak kombinace lichého protonového a neutronového čísla je u stabilních jader vzácná, známe pouze čtyři: \ce{$^{1}_{1}$H}, \ce{$^{6}_{3}$Li}, \ce{$^{10}_{\ 5}$B}, \ce{$^{14}_{\ 7}$N}
\end{itemize}

\textbf{Přeměna $\alpha$} -- uvolňuje se jádro helia, \ce{^4_2He^{2+}}

\ce{$^{226}_{\ 88}$Ra -> $^{222}_{\ 86}$Rn + ^4_2$\alpha$}

\ce{$^{241}_{\ 95}$Am -> ^A_ZX + ^4_2$\alpha$ ($^{237}_{\ 93}$Np)}

\textbf{Přeměna $\beta^-$} -- rozpadem neutronu vzniká elektron

\ce{^1_0n -> ^1_1p + ^0_{-1}$\beta$^-}

\ce{$^{241}_{\ 94}$Pu -> ^A_ZX + ^0_{-1}$\beta$ ($^{241}_{\ 95}$Am)}

\ce{$^{239}_{\ 93}$Np -> ^A_ZX + ^0_{-1}$\beta$ ($^{239}_{\ 94}$Pu)}

\textbf{Přeměna $\beta^+$} -- rozpadem protonu vzniká pozitron

\ce{^1_1p -> ^1_0n + ^0_{1}$\beta$^+}

\ce{$^{23}_{12}$Mg -> ^A_ZX + ^0_{1}$\beta$ ($^{23}_{11}$Na)}

\ce{$^{11}_{\ 6}$C -> ^A_ZX + ^0_{1}$\beta$ ($^{11}_{\ 5}$B)}

\textbf{Elektronový záchyt}

Vnitřní elektron (slupka K) je pohlcen protonem v jádře a vzniká neutron.

\ce{$^{40}_{19}$K -> $^{40}_{18}$Ar}

\textbf{Přeměna $\gamma$}

\ce{$^{60m}_{\ \ 27}$Co -> $^{60}_{27}$Co + $\gamma$}

\subsubsection{Poločas rozpadu}

Statistická veličina, doba za kterou se rozpadne polovina jader v systému.

$\diff{N}{t} = -\lambda N$

$N(t) = N_0e^{-\lambda t}$

$t_{\frac{1}{2}} = \frac{\ln 2}{\lambda} = \tau\ln 2$

\clearpage

\subsection{Elektronový obal}

\begin{itemize}
	\item \textbf{Hlavní kvantové číslo} (\textit{n}) -- číslo slupky: 1, 2, 3, $\ldots$
	\item \textbf{Vedlejší kvantové číslo} (\textit{l}) -- tvar orbitalu: 1 -- \textit{n}-1
	\item \textbf{Magnetické kvantové číslo} (\textit{m}) -- orientaci orbitalu: $<-l;+l>$
	\item \textbf{Spinové kvantové číslo} (\textit{s}) -- popisuje elektron v orbitalu: $\pm\frac{1}{2}$
	\item \textbf{Pauliho princip výlučnosti} - v atomu nemohou existovat dva elektrony, které by měly shodná všechna čtyři kvantová čísla, musí se lišit alespoň spinem, tzn. že do jednoho atomového orbitalu se vejdou maximálně dva elektrony.
	\item \textbf{Výstavbový (Aufbau) princip} - elektrony zaplňují orbitaly od energeticky nejnižších. První jsou zaplňovány volné orbitaly s nejnižším součtem n+l.
	\item Orbitaly jsou zaplňovány v pořadí: 1s, 2s, 2p, 3s, 3p, 4s, 3d, 4p, 5s, 4d, 5p, 6s, 4f, 5d, 6p, 7s, 5f, 6d, 7p
	\item d-orbitaly se zaplňují až po zaplnění s-orbitalu s hlavním kvantovým číslem (n+1), např. 3d orbital se začne plnit až po 4s
\end{itemize}

Zápis elektronové konfigurace vodíku, helia, lithia, sodíku; zkrácený zápis

Tvorba kationtů a aniontů

\begin{tabular}{|l|l||l|l||l|l||l|l|}
	\hline
	H & 1s$^1$ & H$^+$ & 1s$^0$ & B & 1s$^2$ 2s$^2$ 2p$^1$ & B$^{3+}$ & 1s$^2$ 2s$^0$ 2p$^0$ (He) \\\hline
	He & 1s$^2$ & H$^-$ & 1s$^2$ & N & 1s$^2$ 2s$^2$ 2p$^3$ & N$^{3-}$ & 1s$^2$ 2s$^2$ 2p$^6$ (Ne) \\\hline
	Li & 1s$^2$ 2s$^1$ & Li$^+$ & 1s$^2$ & O & 1s$^2$ 2s$^2$ 2p$^4$ & O$^{2-}$ & 1s$^2$ 2s$^2$ 2p$^6$ (Ne) \\\hline
	Be & 1s$^2$ 2s$^2$ & Be$^{2+}$ & 1s$^2$ & F & 1s$^2$ 2s$^2$ 2p$^5$ & F$^{-}$ & 1s$^2$ 2s$^2$ 2p$^6$ (Ne) \\\hline
	Na & [Ne] 3s$^1$ & Na$^{+}$ & [Ne] 3s$^0$ & Si & [Ne] 3s$^2$ 3p$^2$ & Si$^{4-}$ & [Ne] 3s$^2$ 3p$^6$ \\\hline
\end{tabular}

\textbf{Přechodné kovy}: Ca \ce{->} Sc, energie 3d orbitalu; konfigurace Cr a Cu

\begin{itemize}
	\item K: [Ar] 4s$^1$ (3d$^0$)
	\item Ca: [Ar] 4s$^2$ (3d$^0$)
	\item Sc: [Ar] 3d$^1$ 4s$^2$
	\item Ti: [Ar] 3d$^2$ 4s$^2$
	\item Mn: [Ar] 3d$^5$ 4s$^2$
\end{itemize}

Vyšší stabilita zpola a zcela zaplněných d-orbitalů:
\begin{itemize}
	\item Cr: [Ar] 3d$^5$ 4s$^1$
	\item Cu: [Ar] 3d$^{10}$ 4s$^1$
\end{itemize}